\Needspace{12\baselineskip}
\section{Training a Transformer LM}
% ======================================================================
\problem{cross\_entropy}{Implement cross-entropy}{1}

\deliverable{Write a function to compute the cross-entropy loss.}

% Implementation summary:
Subtracting each prediction's maximum
logit avoids exponential overflow without changing the loss. For $n$
predictions and targets $y_b$,
\[
s_{bj}=z_{bj}-\max_k z_{bk},\qquad
\mathcal{L}=\frac1n\sum_{b=1}^{n}
\left[\log\sum_j\exp(s_{bj})-s_{b,y_b}\right].
\]
The target lookup retains a singleton vocabulary axis for matching shapes:
\par\noindent
\begin{minipage}{\linewidth}
\begin{lstlisting}[caption={Stable cross-entropy after shifting logits}]
logits_sum = reduce(torch.exp(shifted), "... vocab_size -> ... 1", "sum")
target_logits = shifted.gather(-1, targets.unsqueeze(-1))
return (torch.log(logits_sum) - target_logits).mean()
\end{lstlisting}
\end{minipage}\par
The provided test passed, including its large-logit case.


\problem{learning\_rate\_tuning}{Tuning the learning rate}{1}

\deliverable{Run the SGD example with learning rates 1e1, 1e2, and 1e3 for 10 iterations. What happens?}

% Answer:
Compared with the slow loss decrease at $\alpha=1$, $\alpha=10$
decreased the loss faster and $\alpha=100$ reached approximately zero within
10 iterations (the first update left the loss essentially unchanged).
At $\alpha=1000$, the loss diverged, increasing from about $24.17$ to
$6.35\times10^{19}$ after 10 updates.

\begin{center}
\begin{tabular}{rrr}
\toprule
Initial learning rate & Initial loss & Loss after 10 updates \\
\midrule
1 & 24.1693 & 19.7481 \\
10 & 24.1693 & 2.8504 \\
100 & 24.1693 & $3.8487\times10^{-24}$ \\
1000 & 24.1693 & $6.3511\times10^{19}$ \\
\bottomrule
\end{tabular}
\end{center}
A local random generator with seed 42 sampled a single $10\times10$ initial
weight matrix. Each run cloned this matrix and created a fresh SGD optimizer,
holding the initialization, objective, precision, and number of updates fixed
while varying only the initial learning rate. All runs used the decaying step
size $\alpha/\sqrt{t+1}$ from the SGD example; losses were recorded initially
and after each of 10 updates.


\problem{adamw}{Implement AdamW}{2}

\deliverable{Implement the AdamW optimizer as a subclass of \texttt{torch.optim.Optimizer}.}

% Implementation summary:
Subclass \texttt{torch.optim.Optimizer},
validate the hyperparameters, and store their defaults in parameter groups.
For each parameter with a gradient, lazily initialize its first and second
moment tensors to zero and maintain a separate update counter starting at 1.
At update $t$, compute
\[
\alpha_t=\alpha\frac{\sqrt{1-\beta_2^t}}{1-\beta_1^t},\qquad
m\leftarrow\beta_1m+(1-\beta_1)g,\qquad
v\leftarrow\beta_2v+(1-\beta_2)g^2.
\]
Apply decoupled weight decay before the moment-adjusted parameter update:
\[
\theta\leftarrow(1-\alpha\lambda)\theta,\qquad
\theta\leftarrow\theta-\alpha_t\frac{m}{\sqrt{v}+\epsilon}.
\]
The update runs under \texttt{torch.no\_grad()}; the decay factor is computed
once per parameter group using its current learning rate:
\par\noindent
\begin{minipage}{\linewidth}
\begin{lstlisting}[caption={In-place AdamW update for one parameter}]
lr_t = lr * sqrt(1 - beta_2**t) / (1 - beta_1**t)
p.mul_(decay_factor)
m.mul_(beta_1).add_(grad, alpha=1 - beta_1)
v.mul_(beta_2).addcmul_(grad, grad, value=1 - beta_2)
denom = v.sqrt().add_(eps)
p.addcdiv_(m, denom, value=-lr_t)
\end{lstlisting}
\end{minipage}\par
Taking a separate square-root tensor preserves the stored second moment.
An optional closure is evaluated with gradients enabled, and its loss is
returned.


\problem{adamw\_accounting}{Resource accounting for training with AdamW}{2}

\Needspace{6\baselineskip}
\textbf{(a)} How much peak memory does running AdamW require? Decompose into parameters, activations, gradients, and optimizer state.

\deliverable{An algebraic expression for each component, as well as the total.}

% Answer:
Let $B$ denote batch size, $T$ context length, $V$ vocabulary
size, $L$ the number of layers, $d$ model width, $h$ the number of attention
heads, and $f=\frac83d$ the feed-forward width. Assume untied input and output
embeddings, bias-free projections, and 4 bytes per FP32 element.

The input embedding and LM head contribute $2Vd$ parameters. Each block has
$4d^2$ attention weights, $3df$ feed-forward weights, and $2d$ RMSNorm weights;
the final RMSNorm adds another $d$. Thus the parameter count is
\[
P=2Vd+L(4d^2+3df+2d)+d
 =2Vd+L(12d^2+2d)+d.
\]
Parameters and gradients each occupy $4P$ bytes. AdamW's two moment tensors
occupy $8P$ bytes; per-parameter counters are neglected.

For activation accounting, retain one output for each listed operator,
without checkpointing or fused attention. The per-block element counts are:
\begin{center}
\begin{tabular}{lr}
\toprule
Operator outputs & Elements per block \\
\midrule
Two RMSNorms & $2BTd$ \\
Q, K, V projections & $3BTd$ \\
Attention scores & $BhT^2$ \\
Softmax probabilities & $BhT^2$ \\
Weighted sum of values & $BTd$ \\
Attention output projection & $BTd$ \\
FFN gate and value projections & $2BTf$ \\
SiLU on the gate branch & $BTf$ \\
Elementwise gate product & $BTf$ \\
FFN output projection & $BTd$ \\
\bottomrule
\end{tabular}
\end{center}
Each head produces its own $T\times T$ score matrix, giving $BhT^2$ rather
than $BT^2$ elements. Summing the table gives
$BT(8d+4f)+2BhT^2$ elements per block. Outside the blocks, count $BTd$ elements
for final RMSNorm, $BTV$ logits, and one $BTV$ log-probability tensor for
cross-entropy. The total activation element count is therefore
\begin{align*}
A(B)&=B\left[L\left(T(8d+4f)+2hT^2\right)+Td+2TV\right]\\
&=B\left[L\left(\frac{56}{3}Td+2hT^2\right)+Td+2TV\right].
\end{align*}
The resulting memory estimate, in bytes, is
\[
\boxed{M(B)=\underbrace{4P}_{\text{parameters}}
             +\underbrace{4P}_{\text{gradients}}
             +\underbrace{8P}_{\text{AdamW state}}
             +\underbrace{4A(B)}_{\text{activations}}.}
\]
This is a simplified retained-output estimate rather than an allocator peak
measurement. It excludes residual additions, RoPE, masks, RMSNorm statistics,
and implementation-specific temporary buffers; actual backward storage and
buffer reuse can change the peak.

\Needspace{6\baselineskip}
\textbf{(b)} Instantiate your answer for GPT-2 XL. What is the maximum batch size within 80GB?

\deliverable{An expression of the form $a \cdot \text{batch\_size} + b$, and a number for the maximum batch size.}

% Answer:
Substitute $V=50{,}257$, $T=1{,}024$, $L=48$, $d=1{,}600$,
and $h=25$. Use the exact algebraic ratio $f=\frac83d$ requested in this
problem, rather than the rounded width of 4,288 in the earlier accounting
problem. This gives $P=1{,}635{,}537{,}600$ and the following memory breakdown,
using $1\,\mathrm{GB}=10^9$ bytes:
\begin{center}
\begin{tabular}{lr}
\toprule
Component & Memory (GB) \\
\midrule
Parameters & 6.542150400 \\
Gradients & 6.542150400 \\
AdamW moments & 13.084300800 \\
Activations & $16.356614144B$ \\
\bottomrule
\end{tabular}
\end{center}
Consequently,
\[
\boxed{M(B)=16.356614144B+26.168601600\quad\mathrm{GB},}
\qquad
B_{\max}=\left\lfloor\frac{80-26.168601600}{16.356614144}\right\rfloor
=\boxed{3}.
\]
The estimate is 75.238444032 GB for $B=3$ and 91.595058176 GB for $B=4$.
If the rounded width $f=4{,}288$ is used instead, then
$P=1{,}640{,}452{,}800$ and
$M(B)=16.373391360B+26.247244800$ GB; the maximum batch size remains 3.

\Needspace{6\baselineskip}
\textbf{(c)} How many FLOPs does running one step of AdamW take?

\deliverable{An algebraic expression, with a brief justification.}

% Answer:
Count scalar additions, subtractions, multiplications,
divisions, and square roots as one operation each. Exclude forward/backward
computation and the scalar hyperparameter arithmetic that does not scale
with the number of parameter elements. A literal expansion of the algorithm
has the following costs per parameter element:
\begin{center}
\begin{tabular}{lr}
\toprule
Operation & FLOPs per element \\
\midrule
Weight decay: $\theta-(\alpha\lambda)\theta$ & 2 \\
First moment: $\beta_1m+(1-\beta_1)g$ & 3 \\
Second moment: $\beta_2v+(1-\beta_2)g^2$ & 4 \\
Parameter update: $\theta-\alpha_t m/(\sqrt{v}+\epsilon)$ & 5 \\
\bottomrule
\end{tabular}
\end{center}
Thus,
\[
F_{\mathrm{AdamW}}\approx14P
=\boxed{22{,}897{,}526{,}400\ \mathrm{FLOPs}}.
\]
The implementation precomputes $1-\alpha\lambda$ once per parameter group,
so its decay operation needs only one multiplication per element. Under the
same counting convention, this reduces its cost to
$13P=21{,}261{,}988{,}800$ FLOPs. The time estimate below uses the literal
$14P$ convention consistently.

\Needspace{6\baselineskip}
\textbf{(d)} Assuming 50\% MFU on an H100, how long would it take to train GPT-2 XL for 400K steps with batch size 1024?

\deliverable{The number of hours training would take, with a brief justification.}

% Answer:
Count the dominant dense matrix multiplications in a forward
pass, with two FLOPs per multiply-accumulate. Per sequence, the attention
projections cost $8Td^2$ per layer, the two attention matrix multiplications
cost $4T^2d$, and the three FFN projections cost $6Tdf$. The LM head adds
$2TdV$. Therefore,
\begin{align*}
F_{\mathrm{forward},1}
 &=L(8Td^2+4T^2d+6Tdf)+2TdV\\
 &=L(24Td^2+4T^2d)+2TdV\\
 &=3{,}506{,}703{,}564{,}800\ \mathrm{FLOPs}.
\end{align*}
This approximation omits elementwise forward operations. For $B=1{,}024$,
assuming backward costs twice forward, one training step costs
\[
F_{\mathrm{step}}=3B F_{\mathrm{forward},1}+14P
=1.0772616248592\times10^{16}\ \mathrm{FLOPs}.
\]
The effective throughput at 50\% MFU is
$0.5\times495\times10^{12}=2.475\times10^{14}$ FLOP/s. Hence,
\[
t_{\mathrm{step}}\approx43.525722217\ \mathrm{s},\qquad
t_{\mathrm{total}}
=\frac{400{,}000F_{\mathrm{step}}}{(2.475\times10^{14})\,3600}
\approx\boxed{4{,}836.19\ \mathrm{hours}}
\approx201.51\ \mathrm{days}.
\]
The stipulated batch size of 1,024 exceeds the single-device memory estimate
in part (b); this is a throughput-based estimate assuming sufficient memory
or gradient accumulation to achieve the effective batch size. Using the
rounded FFN width of 4,288 instead gives approximately 4,850.07 hours.


\problem{learning\_rate\_schedule}{Implement cosine learning rate schedule with warmup}{1}

\deliverable{Write a function that implements the cosine annealing learning rate schedule with warmup.}

% Implementation summary:
Return a scalar learning rate using three
branches: linear warmup for $t<T_w$, cosine decay for $T_w\leq t\leq T_c$,
and the minimum learning rate for $t>T_c$:
\[
\alpha(t)=
\begin{cases}
\alpha_{\max}\,t/T_w, & t<T_w,\\[2pt]
\alpha_{\min}+\frac12(\alpha_{\max}-\alpha_{\min})
\left[1+\cos\left(\pi\frac{t-T_w}{T_c-T_w}\right)\right],
& T_w\leq t\leq T_c,\\[2pt]
\alpha_{\min}, & t>T_c.
\end{cases}
\]
Here $T_c$ is the iteration at which decay ends, so the cosine interval has
length $T_c-T_w$. Its endpoints are $\alpha(T_w)=\alpha_{\max}$ and
$\alpha(T_c)=\alpha_{\min}$, assuming $T_c>T_w$.


\problem{gradient\_clipping}{Implement gradient clipping}{1}

\deliverable{Write a function that implements gradient clipping.}

% Implementation summary:
Collect references to all nonempty parameter
gradients, allowing repeated traversal even when the parameters are supplied
as an iterator. Return immediately if there are no gradients. Compute a single
global norm over all parameter elements,
\[
G=\sqrt{\sum_{p:\,p.\mathrm{grad}\ne\mathrm{None}}
                   \sum_i (p.\mathrm{grad}_i)^2}.
\]
With maximum norm $M$ and $\epsilon=10^{-6}$, apply one shared scale:
\[
g_p\leftarrow\begin{cases}
\dfrac{M}{G+\epsilon}g_p,&G\geq M,\\[4pt]
g_p,&G<M.
\end{cases}
\]
Scaling all gradients in place bounds the combined norm while preserving
the direction of the concatenated gradient vector.

